\scanpage{12}
\margin{cf Artin, Woodshole, pour changement de base}

11) \uline{Supports}
\deleted{\unsure{Désignons|Choisissons} \unsure{donc} \unsure{un|$\forall$} ouvert $U$ de $X$,}

\deleted{\unsure{a}} Faisceau final : $X/X$ lui-même \qquad $X(X')=\{e\}$

Les sous faisceaux du faisceau final sont les ouverts de $X$.

\deleted{Support d'un \unsure{f}}
\deleted{Si $F$ est un \unsure{fai}}

La catégorie des faisceaux $F$ sur $X$ qui sur $U$ sont le faisceau final équivaut à la catégorie des faisceaux sur $Y=X-U$, par le foncteur image directe.

De $=$, faisceaux abéliens sur $X$ nuls sur $U$ $=$ faisceaux abéliens sur $Y$.

Support d'un faisceau abélien

\deleted{$\Gamma_Y(F)\subset\Gamma(F)$ \unsure{Défin.}}

Support d'une section d'un faisceau abélien
\[
\Gamma_Y(F)\subset\Gamma(F)
\]
\[
\underline\Gamma_Y(F)\subset F
\qquad
\underline\Gamma_Y(F)(X')=\Gamma_{Y'}(X',F')
\]

Foncteurs dérivés
\[
\left\{\begin{array}{l}
H^i_Y(X,F)\simeq\operatorname{Ext}^i(X;\mathbb Z_Y,F)\\
\underline H^i_Y(F)\simeq\underline{\operatorname{Ext}}^i(\mathbb Z_Y,F)\\
\phantom{\underline H^i_Y(F)}\text{''}\ R^{i-1}u_*(F|U)\quad\text{si }i\geq2
\end{array}\right.
\]
\margin{si $i=0$ ou $1$, on trouve le noyau et conoyau de $F\longrightarrow u_*(F)$}
\margin{cf SGA 1962 ou Notes Hartshorne 1961}
Suite $\left\{\begin{array}{l}\text{exacte}\\\text{spectrale}\end{array}\right.$ \ldots

12) \uline{Pureté}
$Y\subset X$ réguliers, $\operatorname{codim}_X(Y)=p$ en chaque pt de $Y$
\diagram{\begin{tikzcd}
Y \arrow[dr,dash,"\text{lisse}"'] & & X \arrow[dl,dash,"\text{lisse}"]\\
& k &
\end{tikzcd}}

\margin{\deleted{$\unsure{X}$ \unsure{lisse}}}
\margin{Donner définition de la classe fondamentale locale}
\[
\left\{\begin{array}{lll}
\underline H^i_Y((\mathbb Z/n\mathbb Z)_X)&=0&\text{si }i\ne2p\\
&\simeq\mu_n^{\otimes(-p)}&\text{si }i=2p
\end{array}\right.
\]
\margin{$n$ premier aux car. résiduelles de $Y$}
\mspage{11}
